\problema{Car Pooling}{1094}{src/03-intervalos/1094-car-pooling.cpp}{M}

Manejas un coche con cierta \texttt{capacity} de pasajeros y solo avanzas
hacia adelante, nunca retrocedes. Nos dan \texttt{trips[i] = [numPasajeros,
desde, hasta]}: en el tramo del camino que va de \texttt{desde} a
\texttt{hasta} suben \texttt{numPasajeros} personas, que bajan justo en
\texttt{hasta}. Queremos saber si es posible llevar a todos los pasajeros de
todos los viajes sin exceder la capacidad del coche en ningún punto del
camino.

\paragraph{La idea, con un ejemplo de la vida real.} Piensa en una app de
carpooling universitario: el conductor va publicando en qué kilómetro sube
cada grupo de pasajeros y en qué kilómetro se bajan. En vez de simular el
viaje pasajero por pasajero, basta con anotar, en cada kilómetro del
camino, \emph{cuántos suben} y \emph{cuántos bajan} ahí mismo. Si vas
sumando esos cambios netos conforme avanzas, en todo momento sabes cuántos
pasajeros van sentados en el coche, sin importar el orden en que llegaron
las solicitudes de viaje.

\begin{center}
\ttfamily
\begin{tabular}{l|cccccccc}
posición & 0 & 1 & 2 & 3 & 4 & 5 & 6 & 7\\
\hline
viaje A (+2 en 1, -2 en 5) & & +2 & & & & -2 & &\\
viaje B (+3 en 3, -3 en 7) & & & & +3 & & & & -3\\
ocupación acumulada & 0 & 2 & 2 & 5 & 5 & 3 & 3 & 0\\
\end{tabular}
\end{center}

\noindent Con \texttt{trips = [[2,1,5],[3,3,7]]} y \texttt{capacity = 4}: en
la posición \texttt{3} ya van los \texttt{2} pasajeros del viaje A (que
bajan hasta la posición \texttt{5}) más los \texttt{3} nuevos del viaje B:
$2+3=5 > 4$. No cabe, la respuesta es \texttt{false}.

\paragraph{Cómo lo resuelve el código, paso a paso.} Esta es la técnica de
\emph{arreglo de diferencias} (\emph{difference array}), una variante del
barrido de intervalos.
\begin{enumerate}
  \item Creamos un arreglo \texttt{diff} tan largo como el rango de
        posiciones posible del camino, inicializado en ceros.
  \item Para cada viaje \texttt{(num, desde, hasta)}, sumamos \texttt{num} en
        \texttt{diff[desde]} (aquí suben personas) y restamos \texttt{num}
        en \texttt{diff[hasta]} (aquí bajan).
  \item Recorremos \texttt{diff} de la posición más chica a la más grande,
        acumulando una variable \texttt{ocupacion} que va sumando cada
        valor de \texttt{diff} conforme avanzamos.
  \item En cada posición, si \texttt{ocupacion} supera \texttt{capacity},
        devolvemos \texttt{false} de inmediato: en algún punto del camino el
        coche se llenó de más.
  \item Si terminamos el recorrido sin exceder nunca la capacidad,
        devolvemos \texttt{true}.
\end{enumerate}

\lstinputlisting{src/03-intervalos/1094-car-pooling.cpp}

\complejidad{$O(n + d)$}{$O(d)$}

\paragraph{¿Qué significa esa complejidad?} $n$ es el número de viajes y
$d$ es el rango de posiciones del camino (en LeetCode, acotado a 1000).
Construir el arreglo de diferencias toma $O(n)$ (una suma y una resta por
viaje), y barrerlo entero toma $O(d)$; como $d$ es una constante fija en
este problema, en la práctica el algoritmo es lineal en $n$. El espacio es
$O(d)$ por el arreglo de diferencias, también acotado por una constante.

\paragraph{Consejos para la entrevista (Amazon).} El error más común es
restar en \texttt{hasta - 1} en vez de \texttt{hasta}, pensando que hay que
``incluir'' la última posición ocupada; pero el enunciado dice que los
pasajeros bajan \emph{en} \texttt{hasta}, así que ese punto ya no debe
contarlos. Si el rango de posiciones no viniera acotado por el enunciado
(a diferencia de este problema, donde sí lo está a 1000), la alternativa
general es usar un mapa ordenado (\texttt{std::map}) en vez de un arreglo
fijo, con el mismo principio de sumar en \texttt{desde} y restar en
\texttt{hasta}. Vale la pena mencionar también que este mismo patrón de
``arreglo de diferencias'' resuelve cualquier problema de ``aplica un
cambio a un rango, muchas veces, y pregunta el valor final en cada punto''
sin tener que recorrer cada rango por separado.
